% ------------------------------------------------------------------------
%
\begin{frame}{More proofs by induction}

  Here is again the concatenation of two stacks:
  \begin{align*}
    \fun{cat}(        \el,t) &\xrightarrow{\alpha} t\\
    \fun{cat}(\cons{x}{s},t) &\xrightarrow{\smash{\beta}}
    \cons{x}{\fun{cat}(s,t)}
  \end{align*}
  Here is a functional rewrite system that defines a function
  `\fun{len}' that computes the length of a stack.
  \begin{equation*}
    \begin{array}{r@{\;}l@{\;}l@{}}
      \fun{len}\,\el & \smashedrightarrow{\gamma} & 0;\\
      \fun{len}(\cons{x}{s}) & \smashedrightarrow{\delta}
      & 1 + \fun{len}\, s.
    \end{array}
  \end{equation*}
  So, for example: \(\fun{len}\,[1;2;3] \twoheadrightarrow 1 + (1 + (1
  + (0))) = 3.\)

  \bigskip

  Let us prove the property
  \begin{equation*}
    \pred{LenCat}{s,t} \colon \fun{len}(\fun{cat}(s, t)) \equiv
    \fun{len}\,s + \fun{len}\,t.
  \end{equation*}
  That is: concatenation preserves the length of the stacks.

\end{frame}

% ------------------------------------------------------------------------
%
\begin{frame}{More proofs by induction}

  Note how there are \textbf{two parameters}, \(s\)~and~\(t\), instead
  of one.

  \bigskip

  In this case, we can try to prove the property by induction on one
  parameter, as we did before, and hope that the induction hypothesis
  with the other parameter considered constant suffice. Let us chose
  the first one.

  \bigskip

  The base case is \(\forall t \in S.\pred{LenCat}{\el,t}\), that is
  \[\fun{len}(\fun{cat}(\el,t)) \equiv \fun{len}\,\el +
  \fun{len}\,t.\] It is immediate (the call being rewritten is
  underlined):
  \begin{equation*}
    \fun{len}(\ufun{cat}(\el,t)) \smashedrightarrow{\alpha}
    \fun{len}\,t = 0 + \fun{len}\,t
    \smashedleftarrow{\gamma} \ufun{len}\,\el + \fun{len}\,t.
  \end{equation*}

\end{frame}

% ------------------------------------------------------------------------
%
\begin{frame}{More proofs by induction}

  The general case \(\forall s,t \in S.\pred{LenCat}{s,t} \Rightarrow
  \forall x \in T.\pred{LenCat}{\cons{x}{s}, t}\), where \(S\)~is the
  set of all stacks, and \(T\)~is the set of all terms (\(S \subset
  T\)).

  \bigskip

  Let us assume now \(\forall s,t \in S.\pred{LenCat}{s,t}\) (the
  induction hypothesis), and prove \(\forall x \in
  T.\pred{LenCat}{\cons{x}{s}, t}\).

  \bigskip

  We have:
  \begin{equation*}
    \begin{array}{r@{\;}l@{\;}l@{\qquad}r@{}}
        \fun{len}(\ufun{cat}(\cons{x}{s},t))
      & \smashedrightarrow{\beta}
      & \ufun{len}(\cons{x}{\fun{cat}(s,t)})\\
      & \smashedrightarrow{\delta}
      & 1 + \fun{len}(\fun{cat}(s,t))\\
      & \equiv & 1 + (\fun{len}\,s + \fun{len}\,t)
      & \pred{LenCat}{s,t}\\
      & = & (1 + \fun{len}\,s) + \fun{len}\,t
      & \text{associativity of \((+)\)}\\
      & \smashedleftarrow{\delta}
      & \ufun{len}(\cons{x}{s}) + \fun{len}\,t.
    \end{array}
  \end{equation*}
  We proved \(\forall x \in T.\pred{LenCat}{\cons{x}{s}, t}\).\hfill\(\Box\)

\end{frame}
